%%%%%%%%%%  scan idx 37  |  folio 24  %%%%%%%%%%
spécial~; mais nous n'aurons pas besoin de ce fait).

\subsection*{\S 2. Exemples}

1) Soient $G$ un groupe et $k$ un anneau commutatif~; on considère le groupe $\underline{K}(G)$, quotient du groupe libre engendré par les classes de représentations de $G$ dans des $k$-modules de type fini par le sous-groupe engendré par les éléments de la forme $(\rho\oplus\rho')-(\rho)-(\rho')$ (on note $\oplus$ la somme directe). Le produit tensoriel et les puissances extérieures définissent dans $\underline{K}(G)$ une structure de $\lambda$-anneau. On peut aussi passer au quotient par les éléments de la forme $(\rho)-(\rho')-(\rho'')$ où $\rho$ est une présentation extension de $\rho'$ par $\rho''$, triviale comme extension de $k$-modules. On obtient le groupe $\underline{K}_r(G)$ des classes réduites de représentations de $G$, et l'on constate que la loi de $\lambda$-anneau passe au quotient.

Lorsque $k$ est un corps \emph{algébriquement clos de caractéristique} $0$\footnote{Il est inutile que $k$ soit algébriquement clos, cf.\ (VI 3.3) pour un résultat encore plus général (englobant tous ceux qui seront utilisés dans le présent rapport).}, on peut montrer que $\underline{K}(G)$, et a fortiori $\underline{K}_r(G)$, est un $\lambda$-anneau \emph{spécial}~; il n'en est plus de même si la caractéristique est $\neq 0$, mais cependant si $k$ est algébriquement clos\footnotemark[\value{footnote}], $\underline{K}_r(G)$ est encore un $\lambda$-anneau spécial. Pour le prouver, on est ramené au cas où $G$ est un produit $\Gl(m,k)\times\Gl(n,k)$ et où l'on ne considère que les représentations rationnelles de $G$~; notre assertion résulte de la \emph{détermination explicite} de $\underline{K}_r(G)$ \emph{qui sera faite plus bas} (exemple 3)~; enfin, lorsque $k$ est de caractéristique $0$, la complète réductibilité des représentations rationnelles de $G$ montre que $\underline{K}(G)=\underline{K}_r(G)$.

